%! Setting Up a Test for the Difference of Two Population Means — Unit 7, Lesson 7.8 — Homework (Answer Key)
\documentclass[10pt]{article}
\usepackage{apstats-article}
\usepackage{apstats-key}   % pulls in -boxes; do NOT also load -boxes
\usepackage{grading}
\gradingassignment{apstats-u07-l08-hw}
\gradingcourse{advanced_statistics}
\newcommand{\TallMath}[1]{$\displaystyle #1\rule[-1.4em]{0pt}{3.2em}$}

\begin{document}

\pageheader{Unit 7, Lesson 7.8}{Homework --- Answer Key}

\vspace{0.12in}

\textit{For each problem: (a) name the inference method; (b) define $\mu_1$ and $\mu_2$ and
state $H_0$ and $H_a$; (c) verify conditions. Do \textbf{not} compute a test statistic.}

\begin{enumerate}

\item \textbf{Pain Relief.}

\begin{enumerate}[label=\alph*.]
  \item Method: \ans{two-sample $t$-test for $\mu_1 - \mu_2$}\gradedkey{\gradedblank[id=hw1a, type=short, points=1, answer={{two-sample t-test for mu1 - mu2}}, check=text, display={\boldmath two-sample $t$-test for $\mu_1 - \mu_2$}]{5.3cm}{0.19in}}
  \item $\mu_1 = $ \ans{mean pain score for all patients given the new drug}\gradedkeyline{\gradedblank[id=hw1b1, type=rubric, points=1, display={\bfseries\boldmath mean pain score, new drug}]{\linewidth+6pt}{0.19in}}
        $\mu_2 = $ \ans{mean pain score for all patients given the placebo}\gradedkeyline{\gradedblank[id=hw1b2, type=rubric, points=1, display={\bfseries\boldmath mean pain score, placebo}]{\linewidth+6pt}{0.19in}}

        $H_0$: \ans{$\mu_1 - \mu_2 = 0$} \quad $H_a$: \ans{$\mu_1 - \mu_2 < 0$}\gradedkey{$H_0$:~\gradedblank[id=hw1b3, type=short, points=1, answer={{mu1 - mu2 = 0}}, check=text, display={\boldmath$\mu_1 - \mu_2 = 0$}]{4.3cm}{0.19in}\quad $H_a$:~\gradedblank[id=hw1b4, type=short, points=1, answer={{mu1 - mu2 < 0}}, check=text, display={\boldmath$\mu_1 - \mu_2 < 0$}]{4.3cm}{0.19in}}

        Justify the direction of $H_a$: \ansline{Researchers believe the new drug produces lower mean pain scores, so $\mu_1 < \mu_2$, i.e., $\mu_1 - \mu_2 < 0$.}\gradedkey{\gradedblank[id=hw1b5, type=rubric, points=1, display={\bfseries\boldmath researchers expect the drug to lower mean pain: $\mu_1 < \mu_2$}]{4.6in}{0.19in}}
  \newpage
  \item Conditions:
        \begin{itemize}
          \boxguard
          \item Random: \ansline{Patients were randomly assigned to the two groups --- Random condition met.}\gradedkey{\gradedblank[id=hw1c1, type=rubric, points=1, display={\bfseries\boldmath patients randomly assigned to the groups --- met}]{4.6in}{0.19in}}
          \boxguard
          \item Normal ($n_1 = 24 < 30$, $n_2 = 26 < 30$): \ansline{Both $n < 30$; must check both groups. Both distributions are approximately symmetric with no outliers --- Normal condition met.}\gradedkey{\gradedblank[id=hw1c2, type=rubric, points=1, display={\bfseries\boldmath both $n < 30$; both symmetric, no outliers --- met}]{3.8in}{0.19in}}
        \end{itemize}
\end{enumerate}
\vspace{0.06in}

\item \textbf{Wages.}

\begin{enumerate}[label=\alph*.]
  \item Method: \ans{two-sample $t$-test for $\mu_1 - \mu_2$}\gradedkey{\gradedblank[id=hw2a, type=short, points=1, answer={{two-sample t-test for mu1 - mu2}}, check=text, display={\boldmath two-sample $t$-test for $\mu_1 - \mu_2$}]{5.3cm}{0.19in}}
  \item $H_0$: \ans{$\mu_1 - \mu_2 = 0$} \quad $H_a$: \ans{$\mu_1 - \mu_2 \ne 0$}\gradedkey{$H_0$:~\gradedblank[id=hw2b1, type=short, points=1, answer={{mu1 - mu2 = 0}}, check=text, display={\boldmath$\mu_1 - \mu_2 = 0$}]{4.3cm}{0.19in}\quad $H_a$:~\gradedblank[id=hw2b2, type=short, points=1, answer={{mu1 - mu2 != 0}}, check=text, display={\boldmath$\mu_1 - \mu_2 \ne 0$}]{4.3cm}{0.19in}}
  \newpage
  \item Conditions:
        \begin{itemize}
          \boxguard
          \item Random: \ansline{Both groups are random samples from their respective sectors --- Random condition met.}\gradedkey{\gradedblank[id=hw2c1, type=rubric, points=1, display={\bfseries\boldmath random samples from both sectors --- met}]{4.6in}{0.19in}}
          \boxguard
          \item Normal: \ansline{$n_1 = 50 \ge 30$ and $n_2 = 45 \ge 30$; CLT applies to both --- Normal condition met.}\gradedkey{\gradedblank[id=hw2c2, type=rubric, points=1, display={\bfseries\boldmath $n_1 = 50,\ n_2 = 45 \ge 30$: CLT --- met}]{4.6in}{0.19in}}
        \end{itemize}
\end{enumerate}
\vspace{0.06in}

\item \textbf{Screen Time.}

\begin{enumerate}[label=\alph*.]
  \item Method: \ans{two-sample $t$-test for $\mu_1 - \mu_2$}\gradedkey{\gradedblank[id=hw3a, type=short, points=1, answer={{two-sample t-test for mu1 - mu2}}, check=text, display={\boldmath two-sample $t$-test for $\mu_1 - \mu_2$}]{5.3cm}{0.19in}}
  \item $H_0$: \ans{$\mu_1 - \mu_2 = 0$} \quad $H_a$: \ans{$\mu_1 - \mu_2 < 0$}\gradedkey{$H_0$:~\gradedblank[id=hw3b1, type=short, points=1, answer={{mu1 - mu2 = 0}}, check=text, display={\boldmath$\mu_1 - \mu_2 = 0$}]{4.3cm}{0.19in}\quad $H_a$:~\gradedblank[id=hw3b2, type=short, points=1, answer={{mu1 - mu2 < 0}}, check=text, display={\boldmath$\mu_1 - \mu_2 < 0$}]{4.3cm}{0.19in}}
  \item Conditions:
        \begin{itemize}
          \boxguard
          \item Random: \ansline{Both groups were randomly selected --- Random condition met.}\gradedkey{\gradedblank[id=hw3c1, type=rubric, points=1, display={\bfseries\boldmath both groups randomly selected --- met}]{4.6in}{0.19in}}
          \newpage
          \boxguard
          \item Normal (both $n < 30$): \ansline{$n_1 = 18$ and $n_2 = 21$ are both $< 30$; must check both. Group 1: mild left skew with no outliers --- acceptable. Group 2: approximately symmetric --- acceptable. Normal condition met for both.}\gradedkey{\gradedblank[id=hw3c2, type=rubric, points=1, display={\bfseries\boldmath both $n < 30$; mild skew / symmetric, no outliers --- met}]{4.6in}{0.19in}}

                Can the test proceed? \ans{Yes} \quad Explain: \ansline{Mild skew without outliers is generally not a severe enough violation to prevent inference when $n < 30$.}\gradedkey{\gradedblank[id=hw3c3, type=short, points=1, answer=yes, check=text, display={Yes}]{1.3cm}{0.19in}}
                \gradedkeyline{\gradedblank[id=hw3c4, type=rubric, points=1, display={\bfseries\boldmath mild skew, no outliers: OK}]{\linewidth+6pt}{0.19in}}
        \end{itemize}
\end{enumerate}
\vspace{0.06in}

\newpage
\item \textbf{AP-Style.}

\begin{enumerate}[label=(\Alph*)]
  \item $n_1 + n_2 = 30$ does not satisfy the condition --- the CLT requires each group $\ge 30$.
  \item \textcolor{keyred}{\textbf{$\leftarrow$ correct}} Moderate skewness \emph{without} outliers can be acceptable for $n < 30$; however, \emph{one outlier per group} is a concern.
  \item This is too strict --- moderate skewness alone doesn't violate the condition.
  \item The condition applies to \emph{both} samples, not just one.
\end{enumerate}
\vspace{0.04in}
Answer: \ans{C is the best answer given the choices, but note the outliers make this borderline.}\gradedkey{\textit{Letter:}~\gradedblank[id=hw4a, type=choice, points=1, answer=C]{0.4in}{0.2in}}

Brief justification: \ansline{Both $n < 30$; must check both. Moderate skewness can be tolerated, but outliers in both groups raise concern. Proceed with caution.}\gradedkeyline{\gradedblank[id=hw4b, type=rubric, points=1, display={\bfseries\boldmath both $n < 30$ with an outlier in each group}]{\linewidth+6pt}{0.19in}}

\begin{teachernote}
Question 4 is intentionally nuanced. The AP exam expects students to note that outliers
with small samples are a genuine concern. Accept C with a discussion of the outliers.
\end{teachernote}

\end{enumerate}

\begin{extensionbox}
A student argues: ``Since the null hypothesis is always $\mu_1 - \mu_2 = 0$, a two-sample
$t$-test can never detect a specific numerical difference like $\mu_1 - \mu_2 = 5$.''
Is this correct? Explain when a researcher might hypothesize a non-zero difference and
how the null hypothesis would be written in that case.
\ansline{The student is incorrect. While the standard null hypothesis is $H_0: \mu_1 - \mu_2 = 0$, a researcher can set any hypothesized difference: e.g., $H_0: \mu_1 - \mu_2 = 5$. In that case the test statistic uses 5 as the null value instead of 0. Such non-zero nulls arise when a researcher wants to test whether a known or practically meaningful gap exists between populations.}
\ansline{}
\ansline{}
\end{extensionbox}

\end{document}
